% GENERATED FILE. Edit canonical lessons, then run npm run generate:books.
% canonical-source-hash: fnv1a64:3c81ae0ead3b08ea
% canonical-lessons: TE-C223-vyaktigata, TE-C223-praivetu, TE-C223-panikirani, TE-C223-cindaravandara, TE-C223-moratu

\chapter{Personal, Private, Useless, Messy, Rude}
\label{ch:te-personal-private-useless-messy-rude}
\hlchaptermodality{223}

\begin{chapteropening}
\textbf{By the end of this chapter:} \emph{I can say personal, private, useless, messy and rude.}

Complete the last lesson of chapter 223: I can say personal, private, useless, messy and rude.
\end{chapteropening}

\section[\texorpdfstring{\te{వ్యక్తిగత} (vyaktigata)}{వ్యక్తిగత (vyaktigata)}]{\te{వ్యక్తిగత} (vyaktigata) — personal}

\label{lesson:TE-C223-vyaktigata}



\begin{quote}
Before the new one: say the Telugu for different, then the Telugu for fair.
\end{quote}

\subsection*{You'll want to know: \te{వ్యక్తిగత}}
\textbf{\te{వ్యక్తిగత}} — \emph{vyaktigata} — \textquotedblleft{}personal\textquotedblright{}.

\begin{grammarlens}[title={What you've built}]
One more word for this chapter.
\end{grammarlens}

\subsection*{Your turn}
\begin{itemize}
\raggedright
  \item \emph{Say it:} \emph{vyaktigata}
  \item \emph{Say it:} \emph{vyaktigata}, once more
  \item \emph{Say it:} say \emph{vyaktigata}
  \item \emph{Recall:} say the Telugu for different, then the Telugu for fair, then say \emph{vyaktigata} again
\end{itemize}

\begin{tcolorbox}[breakable,colback=teal!4,colframe=teal!35!black,title={Before you move on}]
What does \te{వ్యక్తిగత} mean? (\textquotedblleft{}Personal\textquotedblright{}.) Say it once more.
\end{tcolorbox}
\section[\texorpdfstring{\te{ప్రైవేటు} (praivēṭu)}{ప్రైవేటు (praivēṭu)}]{\te{ప్రైవేటు} (praivēṭu) — private}

\label{lesson:TE-C223-praivetu}



\begin{quote}
Before the new one: say the Telugu for fair, then the Telugu for personal.
\end{quote}

\subsection*{You'll want to know: \te{ప్రైవేటు}}
\textbf{\te{ప్రైవేటు}} — \emph{praivēṭu} — \textquotedblleft{}private\textquotedblright{}.

\begin{grammarlens}[title={What you've built}]
One more word for this chapter.
\end{grammarlens}

\subsection*{Your turn}
\begin{itemize}
\raggedright
  \item \emph{Say it:} \emph{praivēṭu}
  \item \emph{Say it:} \emph{praivēṭu}, once more
  \item \emph{Say it:} say \emph{praivēṭu}
  \item \emph{Recall:} say the Telugu for fair, then the Telugu for personal, then say \emph{praivēṭu} again
\end{itemize}

\begin{tcolorbox}[breakable,colback=teal!4,colframe=teal!35!black,title={Before you move on}]
What does \te{ప్రైవేటు} mean? (\textquotedblleft{}Private\textquotedblright{}.) Say it once more.
\end{tcolorbox}
\section[\texorpdfstring{\te{పనికిరాని} (panikirāni)}{పనికిరాని (panikirāni)}]{\te{పనికిరాని} (panikirāni) — useless}

\label{lesson:TE-C223-panikirani}



\begin{quote}
Before the new one: say the Telugu for personal, then the Telugu for private.
\end{quote}

\subsection*{You'll want to know: \te{పనికిరాని}}
\textbf{\te{పనికిరాని}} — \emph{panikirāni} — \textquotedblleft{}useless\textquotedblright{}.

\begin{grammarlens}[title={What you've built}]
One more word for this chapter.
\end{grammarlens}

\subsection*{Your turn}
\begin{itemize}
\raggedright
  \item \emph{Say it:} \emph{panikirāni}
  \item \emph{Say it:} \emph{panikirāni}, once more
  \item \emph{Say it:} say \emph{panikirāni}
  \item \emph{Recall:} say the Telugu for personal, then the Telugu for private, then say \emph{panikirāni} again
\end{itemize}

\begin{tcolorbox}[breakable,colback=teal!4,colframe=teal!35!black,title={Before you move on}]
What does \te{పనికిరాని} mean? (\textquotedblleft{}Useless\textquotedblright{}.) Say it once more.
\end{tcolorbox}
\section[\texorpdfstring{\te{చిందరవందర} (cindaravandara)}{చిందరవందర (cindaravandara)}]{\te{చిందరవందర} (cindaravandara) — messy, in disarray}

\label{lesson:TE-C223-cindaravandara}



\begin{quote}
Before the new one: say the Telugu for private, then the Telugu for useless.
\end{quote}

\subsection*{You'll want to know: \te{చిందరవందర}}
\textbf{\te{చిందరవందర}} — \emph{cindaravandara} — \textquotedblleft{}messy, in disarray\textquotedblright{}.

\begin{grammarlens}[title={What you've built}]
One more word for this chapter.
\end{grammarlens}

\subsection*{Your turn}
\begin{itemize}
\raggedright
  \item \emph{Say it:} \emph{cindaravandara}
  \item \emph{Say it:} \emph{cindaravandara}, once more
  \item \emph{Say it:} say \emph{cindaravandara}
  \item \emph{Recall:} say the Telugu for private, then the Telugu for useless, then say \emph{cindaravandara} again
\end{itemize}

\begin{tcolorbox}[breakable,colback=teal!4,colframe=teal!35!black,title={Before you move on}]
What does \te{చిందరవందర} mean? (\textquotedblleft{}Messy, in disarray\textquotedblright{}.) Say it once more.
\end{tcolorbox}
\section[\texorpdfstring{\te{మొరటు} (moraṭu)}{మొరటు (moraṭu)}]{\te{మొరటు} (moraṭu) — rude, coarse}

\label{lesson:TE-C223-moratu}



\begin{quote}
Before the new one: say the Telugu for useless, then the Telugu for messy.
\end{quote}

\subsection*{You'll want to know: \te{మొరటు}}
\textbf{\te{మొరటు}} — \emph{moraṭu} — \textquotedblleft{}rude, coarse\textquotedblright{}.

\begin{grammarlens}[title={What you've built}]
One more word for this chapter.
\end{grammarlens}

\subsection*{Your turn}
\begin{itemize}
\raggedright
  \item \emph{Say it:} \emph{moraṭu}
  \item \emph{Say it:} \emph{moraṭu}, once more
  \item \emph{Say it:} say \emph{moraṭu}
  \item \emph{Recall:} say the Telugu for useless, then the Telugu for messy, then say \emph{moraṭu} again
\end{itemize}

\begin{tcolorbox}[breakable,colback=teal!4,colframe=teal!35!black,title={Before you move on}]
What does \te{మొరటు} mean? (\textquotedblleft{}Rude, coarse\textquotedblright{}.) Say it once more.
\end{tcolorbox}
